\documentclass{beamer}

% Theme selection (you can change to Madrid, Berlin, Warsaw, etc.)
\usetheme{CambridgeUS}

\usepackage{fontspec}        % 文本字体
\usepackage{unicode-math}    % 数学字体

% Presentation metadata
\title{Fast Fourier Transform}
\author{Wenyue Zhang}
%\institute{Department of Mathematics, Your University}
\date{\today}

\begin{document}

% Title slide
\begin{frame}
  \titlepage
\end{frame}

% Table of contents (optional)
\begin{frame}{Outline}
  \tableofcontents
\end{frame}

%------------------------------------------------
\section{Fourier Analysis on $\mathbb Z(N)$}

\begin{frame}{The Group $\mathbb Z(N)$}
    Two integers $x$ and $y$ are congruent modulo $N$ if $x - y$ is divisible by $N$.

    \pause

    It can be easily seen that this defines an equivalence relation on $\mathbb Z$.

    \pause

    Define $R(x)$ the equivalence class of $x$, and define $$R(x) + R(y) = R(x + y).$$This turns the set of equivalence classes into an abelian group called the 
    group of integers modulo $N$ denoted by $\mathbb Z(N)$ or $\mathbb Z / N \mathbb Z$.
\end{frame}

\begin{frame}{Root of Unity}
    A complex number $z$ is an $N$th root of unity of $z^N = 1$. The set of $N$th roots of unity is
    $$
    \{1, \text{e}^{2 \pi \text{i}/N}, \text{e}^{2(2 \pi \text{i}/N)}, \cdots, \text{e}^{(n - 1)(2 \pi \text{i}/N)}\}
    $$
    It can be seen that the set forms a group under complex multiplication. Moreover, it is essentially the same as $\mathbb Z(N)$.

    \pause

    We will use $\mathbb Z(N)$ to denote either of these two groups.
\end{frame}

\begin{frame}{Fourier Coefficient}
    Assume $f : [0, 2 \pi) \to \mathbb C$ is a $2 \pi$-periodic function. The Fourier coefficient of $f$ is defined by
    $$
    a_n = \frac{1}{2 \pi}\int_0^{2 \pi} f(x) \text{e}^{-\text{i} n x} \text{d}x.
    $$
    \pause
    Under certain conditions, we have
    $$
    f(x) = \sum_{n = - \infty}^{\infty} a_n \text{e}^{\text{i}nx}
    $$
\end{frame}

\begin{frame}{Fourier Coefficient}
    Let $V$ be the set of all continuous functions on $[0, 2 \pi)$. It is clear that $V$ is a vector space under addition and scaler multiplication.
    \pause
    If we define the inner product $\langle f, g \rangle$ by
    $$
    \langle f, g \rangle = \frac{1}{2 \pi}\int_0^{2 \pi} f \overline g,
    $$
    then one can verify that $\{e_n\}$ form an orthonormal basis of $V$, where $e_n(x) = \text{e}^{\text{i}nx}$.
\end{frame}

\begin{frame}{Fourier Coefficient}
    Now let consider all complex-valued functions on $\mathbb Z(N)$. Similarly denote this vector space by $V$. We define the inner product by
    $$
    \langle f, g \rangle = \sum_{n = 0}^{N - 1} f(n) \overline{g(n)}.
    $$
    \pause
    If we define $e_l(k) = \text{e}^{2 \pi \text{i}lk}$ for $l = 0, 1, \cdots, N - 1$ and $k = 0, 1, \cdots, N - 1$, then it can be verified that 
    $\{e_0, \cdots, e_{N - 1}\}$ is an orthogonal basis.
\end{frame}

\begin{frame}{Fourier Coefficient}
    By results from linear algebra, we can write $f$ as
    $$
    f = \sum_{n = 0}^{N - 1} a_n e_n,
    $$
    where
    $$
    a_n = \frac{\langle f, e_n \rangle}{\langle e_n, e_n \rangle} = \frac{1}{N}\sum_{k = 0}^{N - 1} f(k) \text{e}^{- 2 \pi \text{i} kn / N}.
    $$
    \pause
    In addition,
    $$
    f(k) = \sum_{n = 0}^{N - 1} a_n \text{e}^{2 \pi \text{i} nk/N}
    $$
\end{frame}

%------------------------------------------------
\section{Fast Fourier Transform}

\begin{frame}{Discrete Fourier Transform}
    The process to find $a_0, a_1, \cdots, a_{N - 1}$ is called the discrete Fourier transform.

    \pause

    We can use two loops to compute them, with a time complexity of $O(N^2)$.

    \pause

    We now present the fast Fourier transform, an algorithm that improves the bound $O(N^2)$ to $O(N \log N)$.
\end{frame}

% Formula example
\begin{frame}{Fast Fourier Transform}
    The algorithm is based on the technique of divide-and-conquer.

    \pause

    First we restrict ourselves to the case where $N = 2^n$. After that it is easy to extend the method to all integers by adding $0$ term.

    \pause

    Let $\#(M)$ denote the minimum number of operations needed to calculate all the Fourier coefficients. The key is contained in the following recursion step.
\end{frame}

\begin{frame}{A Lemma}
    If we are given $\omega_{2M} = \text{e}^{-2 \pi \text{i} / (2M)}$, then
    $$
    \#(2M) \le 2\#(M) + 8M
    $$
    \pause
    Proof: The calculation of $\omega_{2M}, \cdots, \omega_{2M}^{2M}$ requires no more than $2M$ operations. Note that in particular we get $\omega_M = \omega_{2M}^2$. 
    The main idea is that for any given function $f$ on $\mathbb Z(M)$, we consider two functions $f_0$ and $f_1$ on $\mathbb Z(M)$ defined by
    $$
    f_0(r) = f(2r) \text{ and } f_1(r) = f(2r + 1).
    $$
\end{frame}

\begin{frame}{A Lemma}
    We assume that it is possible to calculate the Fourier coefficient of $F_0$ and $F_1$ in no more than $\#(M)$ operations each. If we denote the Fourier coefficients corresponding to 
    the groups $\mathbb Z(2M)$ and $\mathbb Z(M)$ by $a_{k}^{2M}$ and $a_{k}^M$, then we have
    $$
    \begin{aligned}
        a_k^{2M}(f) &= \frac{1}{2M} \sum_{r = 0}^{2M - 1} f(r) \omega_{2M}^{kr} \\
        &= \frac{1}{2}\left(\frac{1}{M}\sum_{r = 0}^{M - 1}f_0(r) \omega_M^{kr} + \frac{1}{M}\sum_{r = 0}^{M - 1} f_1(r) \omega_M^{km} \omega_{2M}^k\right).
    \end{aligned}
    $$
\end{frame}

\begin{frame}{A Lemma}
    Therefore
    $$
    a_k^{2M}(f) = \frac{1}{2} \left(a_k^M(f_0) + a_k^M(f_1) \omega_{2M}^k \right).
    $$
    As a result, knowing $a_k^M(f_0), a_k^M(f_1)$ and $\omega_{2M}^k$, we see that each $a_k^{2M}(f)$ can be computed using no more than three operations. So
    $$
    \#(2M) \le 2M + 2 \#(M) + 3 \times 2M = 2 \#(M) + 8M.
    $$
\end{frame}

\begin{frame}{Fast Fourier Transform}
    Then we can prove the bound by induction. The initial step for $n = 1$ is easy, since $N = 2$ and the two Fourier coefficients can be computed easily. Assume $\#(N) \le 4 \cdot 2^nn$, then by the lemma
    $$
    \#(2N) \le 2 \#(N) + 8N = 2 \cdot 4 \cdot 2^nn + 8 \cdot 2^n = 4\cdot2^{n + 1}(n + 1).
    $$
    Hence we conclude that $\#(N) = O(N \log N)$.
\end{frame}

% Image example (place image.png in the same folder)
% \begin{frame}{Illustration}
%   \centering
%   \includegraphics[width=0.6\textwidth]{image.png}
% \end{frame}

%------------------------------------------------
\section{Polynomial Multiplication}

\begin{frame}{Polynomial Multiplication}
    We present one application of fast Fourier transform.

    \vspace{1em}

    \pause

    Given polynomials $F$ and $G$, compute their product $F \cdot G$. Assume the degree of $F$ and $G$ is both $n - 1$, otherwise we can add additional $0$ term to satisfy the condition.

    \vspace{1em}

    \pause

    If we directly compute every term in $F \cdot G$, the time complexity would be $O(n^2)$.
\end{frame}

\begin{frame}{Process}
    First let extend $F$ and $G$ to a polynomial of degree of $2n - 2$, since the degree of $F \cdot G$ is $2n - 2$.
    
    \pause

    Now let us get the Fourier coefficients of $F$ and $G$ by fast Fourier transform, thus $F(k) = \sum a_n \text{e}^{2 \pi\text{i}nk/N}$ and $G(k) = \sum b_n \text{e}^{2 \pi \text{i}nk/N}$. 
    Assume the Fourier coefficients of $F \cdot G$ are $\{c_n\}$, then $c_n = a_nb_n$(why?).

    \pause

    Since we know the Fourier coefficients of $F \cdot G$, we surely know $F \cdot G$.
    
\end{frame}

\begin{frame}{Time Complexity}
    The time complexity for getting $\{a_n\}$ and $\{b_n\}$ is $O(n \log n)$ by fast Fourier transform.

    \pause

    It takes $O(n)$ time to compute $\{c_n\}$.

    \pause

    The process to get $F \cdot G$ from $\{c_n\}$ is called the inverse fast Fourier transform. It can be done by modifying the algorithm for getting Fourier coefficient and the time complexity is thus $O(n \log n)$.

    \pause

    The overall time complexity is $O(n \log n)$.
\end{frame}

%------------------------------------------------
\section{Conclusion}

% Acknowledgements
\begin{frame}{Acknowledgements}
    I am deeply grateful to Frank Cole for his help and guidance, DRP coordinators for organizing this and Elias M. Stein \& Rami Shakarchi for their great textbook.

    \pause

    \vspace{1em}

    Thanks for listening.
\end{frame}

\end{document}
